\documentclass[11pt]{book}

\input{.house-style/preamble.tex}

\hypersetup{
  pdfauthor={Brett Reynolds},
  pdftitle={Book Proposal: Words That Won't Hold Still}
}

\renewcommand{\thesection}{\arabic{section}}
\renewcommand{\thesubsection}{\arabic{section}.\arabic{subsection}}

\newcommand{\labelsc}[1]{{\mdseries\textsc{#1}}}
\newcommand*{\orcid}{\orcidlink{0000-0003-0073-7195}}

\begin{document}
\pagestyle{fancy}

\chapter*{Book Proposal}
\markboth{Book Proposal}{Book Proposal}
\thispagestyle{plain}
\begin{center}
{\LARGE \emph{Words That Won't Hold Still} \par}
\vspace{0.75em}
{\large \textsc{Cambridge Studies in Linguistics} \par}
\vspace{0.75em}
\labelsc{Author:} Brett Reynolds\,\orcid \par
\labelsc{Affiliation:} Humber Polytechnic \& University of Toronto \par
\labelsc{Manuscript status:} Complete (494 pages, 18 chapters)
\end{center}

\hrule

\section{Overview}

This book begins from a simple but surprisingly neglected problem: linguistics relies on categories whose ontological status it rarely examines. Every day, linguists ask whether a form is a noun, whether a construction is grammatical, whether two sounds belong to the same phoneme, whether a category such as \term{evidential} or \term{resultative} is cross-linguistically real, or whether a borderline case falls inside or outside a word class. Those questions are ordinarily treated as if they presuppose a familiar answer: categories are determinate kinds with discoverable boundaries, and the analyst's task is to locate those boundaries correctly. But the empirical record of linguistics has never sat comfortably with that assumption. Some of the field's most central categories are robust enough to support generalization and explanation, but messy enough to resist essentialist definition.

The usual options have been unsatisfying. On one side is essentialism: categories are defined by necessary and sufficient conditions, membership is binary, and apparent counterexamples are either misdescribed or marginal. This remains the default picture in much linguistic argument, even when it's not stated explicitly. But essentialist definitions fail across domains. Countability doesn't reduce to a single lexical essence. Definiteness is coherent enough to matter but too heterogeneous to be captured by any one simple criterion. Lexical categories, pronouns, reciprocals, constructions, and phonological contrasts show the same broader pattern: asymmetrical clustering and diagnostic conflict, especially at points of overlap. Essentialism preserves determinacy at the cost of descriptive accuracy.

On the other side is the retreat to prototypes, gradients, or nominalist convenience. These approaches rightly register the empirical pattern of probabilistic distributions, family resemblances, and instability near category boundaries. But they leave a central question unanswered: why do these patterns remain stable enough to underwrite linguistic explanation? If categories are only loose statistical similarities, why do they endure across speakers, acquisition, interaction, and transmission? Why do they support induction? Why do speakers converge on them? Prototype descriptions often register gradience, but they don't explain the persistence and projectibility that linguistic theory depends on.

This book argues that linguistics has lacked a fully articulated third option. Drawing on Richard Boyd's theory of homeostatic property-cluster kinds, Ruth Millikan's work on copied kinds, and related philosophy of science, it proposes that many linguistic categories are real kinds without essences. The book's central claim can be stated as a slogan: \emph{a linguistic category is a profile of co-occurring properties, stabilized by mechanisms, projectible relative to purposes}. On this view, no single property need be necessary; categories may be internally graded; and apparent gradience near boundaries needn't imply gradient membership. What makes the category real is not a hidden essence but a stable clustering of properties maintained by causal processes such as acquisition, entrenchment, interactive alignment, social coordination, and iterated transmission.

The framework is imported from philosophy of biology, where homeostatic property-cluster theory helped resolve the species problem, but the transfer to linguistics isn't a loose analogy. Linguistic categories are also population-level phenomena: they are distributed across speakers, instantiated in usage, acquired by learners, aligned in interaction, and transmitted across generations. These mechanisms explain why categories can tolerate peripheral variation while maintaining a coherent core. They also explain how categories fray, split, merge, or fail. A category such as the English definite article remains coherent not because every use satisfies one invariant semantic condition, but because a cluster of formal and functional properties is continually reproduced and regulated by identifiable mechanisms.

That reconceptualization solves a longstanding explanatory problem. It preserves what linguists need from realism: categories aren't arbitrary labels, and they can ground inductions. But it rejects the assumption that realism requires essences. Borderline cases stop being mere embarrassment and become data about the limits of a category's stabilizing mechanisms. Gradient acceptability judgments, corpus frequency distributions, and experimental probes of category membership become principled tools rather than ad hoc supplements.

A further major contribution is the book's account of field-relative projectibility. Different subfields often ask different explanatory questions of the same material and then treat the resulting disagreement as if one side has to have misdescribed the facts. The manuscript argues instead that overlapping but distinct homeostatic kinds may coexist within the same causal system. \term{proper noun} may project distributional facts for a syntactician; \term{proper name} may project referential facts for a semanticist. These don't need to be rival definitions of one confused entity; instead, they can be distinct but overlapping profiles, each real relative to a field's explanatory purposes. This idea has wide consequences for debates about lexical categories, constructions, grammatical meaning, and typological comparison.

The book doesn't stop at a general philosophical proposal. It works the framework through a broad range of substantive linguistic domains. Countability is recast as a cover term for two tightly coupled categories: a semantic individuation cluster and a morphosyntactic count cluster. Definiteness is split into two conflated HPCs, separating semantic definiteness from a formal cluster and motivating the category of \enquote{deitality}. Lexical categories are re-examined in terms of mechanism density, while gender, sign language anaphora, phonological contrast, constructions, register, and grammaticality itself provide further test cases. Agent-based models show how stable clustering can arise and how it behaves under stress. The final chapter doesn't protect the theory from failure; it assembles predictions, defeat conditions, and explicit falsifiers, turning the framework into a research program that can succeed or fail empirically. Because the manuscript is complete at 494 pages and 18 chapters, reviewers are assessing not a promising sketch but a finished theoretical architecture.

Getting the metaphysics right has consequences across the discipline. Methodologically, it reconceives diagnostics, boundary phenomena, and gradient data: the question is no longer \enquote{How do we explain away the cases that don't fit?} but \enquote{What do the edge cases show about the mechanisms that hold the cluster together?} Theoretically, it dissolves the stale opposition between discreteness and gradience: categories can yield gradient judgments near boundaries while still supporting discrete analytical choices because stability is mechanistic rather than essentialist. More broadly, the book creates a meeting point for traditions that have each seen part of the truth. Formal and descriptive linguistics are right that grammatical categories support strong generalizations. Usage-based and constructional work are right that those categories are frequency-sensitive, learned, and historically contingent. Typologists are right to worry that some cross-linguistic categories are analysts' taxonomies rather than genuine kinds. Sociolinguists are right that inter-speaker alignment and community norms are part of grammatical stability. The HPC framework explains how these insights belong together.

Cambridge Studies in Linguistics is the right home for this book because the series is explicitly dedicated to scholarship that makes a substantive contribution to general, theoretical, and descriptive linguistics. This manuscript does so by addressing a foundational question that cuts across all three: what sort of entity a linguistic category is, and what makes it projectible enough to support explanation. The book is broad without being diffuse, philosophically serious without ceasing to be linguistic, and empirically grounded across multiple domains of analysis. It is a book about what linguistic categories are and what makes them projectible enough to support explanation.

\section{Chapter Summary}

\subsection*{Part I: The Problem}

\labelsc{Chapter 1:} Words That Won't Hold Still -- Linguists assume categories have determinate boundaries; this assumption is rarely examined and systematically misleading. The impasse between essentialism (fails on data) and prototype theory (fails on explanation) leaves the field without adequate foundations.

\labelsc{Chapter 2:} Essentialism and its Discontents -- The default essentialist view fails reliably across traditions and domains. The \emph{CGEL} verbless-clause debate as microcosm. Biology faced an analogous problem with species, and HPC theory became one influential response.

\labelsc{Chapter 3:} What We Haven't Been Asking -- Framework-free approaches to typology (Haspelmath, Croft) avoid essentialism by replacing language-specific categories with comparative concepts, but they leave the stability of language-specific categories unexplained. Usage-based, constructionist, and sociolinguistic traditions have \enquote{the right instincts}; the missing piece is mechanistic metaphysics.

\subsection*{Part II: The Fix}

\labelsc{Chapter 4:} Categories Without Essences -- Boyd's HPC framework and Millikan's copied kinds applied to grammar. The species-grammar parallel is structural, not metaphorical. Five stabilizing mechanisms sketched. Engagement with Ambridge's radical exemplarism, Khalidi's simple causal critique, Slater's stability taxonomy, and Magnus's fundamentality objection.

\labelsc{Chapter 5:} Discrete from Continuous -- How continuous substrates produce discrete categories. Phase transitions as model. Hyperreal formalization. Three worked examples: \emph{fun} drift (English adjective), English reciprocals, Shilluk number marking.

\labelsc{Chapter 6:} The Stabilizers -- Mechanism inventory along timescale and locus dimensions. Case study 1: quotatives cross-linguistically. Case study 2: independent relative \emph{whose}. The chapter argues that mechanisms are themselves categories maintained by further mechanisms -- a recursive architecture the rest of the book exploits.

\labelsc{Chapter 7:} Projectibility and the Good Bet -- What mechanisms buy: reliable inductive inference. The book's central slogan introduced. Polish aspect as case study. Peircean apparatus (interpretant, habit, would-be). Field-relative projectibility with worked examples. Degrees of projectibility.

\labelsc{Chapter 8:} Failure Modes -- Thin clustering, fat clustering, and negative diagnosis. The inflation problem in typology. Two-diagnostic stress test. Seven-step audit template. Coupling continuum from hard (phonemes) to loose (grammatical categories) to composite (constructions).

\subsection*{Part III: Categories Reconsidered}

\labelsc{Chapter 9:} Countability -- Cover term for two tightly coupled categories: semantic individuation cluster and morphosyntactic count cluster. Implicational hierarchy. Object-mass nouns as boundary cases. Diachronic evidence (\emph{pea}, \emph{data}, \emph{cattle}). Agent-based simulation with stress tests. Cross-linguistic variation.

\labelsc{Chapter 10:} Definiteness and Deitality -- Two conflated HPCs: semantic definiteness and article-system form cluster. Millikan's proper vs derived function. Weak and generic definites as derived functions. \enquote{Deitality} proposed for form cluster. Greek and Japanese comparisons.

\labelsc{Chapter 11:} Lexical Categories -- Gradient of mechanism density across word classes. Noun/verb as high-density structural core, adjective as variable density, adverb as thin clustering, pronouns as braided sub-clusters. Diachronic and acquisitional evidence.

\labelsc{Chapter 12:} Pro-form Gender -- English gender as system-wide personhood partition. Maintenance spectrum running from English (transparent semantic basis) through German (hybrid) to French (morphologically entrenched). Chain-coherence phenomena (pronoun-antecedent agreement as maintenance check).

\labelsc{Chapter 13:} Sign Language Anaphor -- Modality independence test. Two-cluster architecture survives in sign: functional cluster (reciprocal semantics, valency reduction) + formal cluster (spatial loci, reduplication). Cross-sign variation (ASL, BSL, village sign languages).

\labelsc{Chapter 14:} The Category Zipper -- Scale-invariance from phonemes to constructions. Four coupling tiers. Designed-to-fail controls. Negative cases (Russian zero-marked definiteness fails HPC audit).

\subsection*{Part IV: What Changes}

\labelsc{Chapter 15:} Grammaticality Itself -- Grammaticality as maintained structural coupling, not semantic compositionality. Ungrammaticality as noisy detector (illusions, satiation, individual differences all follow from mechanistic account).

\labelsc{Chapter 16:} The Social Stabilization of Kinds -- Inter-speaker convergence mechanisms. Categories as coordination equilibria. Dialect and register as two mechanistic dimensions. Belfast social network study as empirical case.

\labelsc{Chapter 17:} How Categories Travel -- Cross-domain extension as predicted capacity of HPCs. Boyd's theory-constitutive metaphor applied to linguistic category extension. Worked example: how a tracking category in canine cognition extends from olfactory to visual to abstract domains. Engagement with cognitive linguistics on its own terms via structured rules of dialogue.

\labelsc{Chapter 18:} The Gauntlet -- Predictions, stress tests, and defeat conditions. Synthesis ledger compares HPC predictions against alternatives across five empirical domains. Two novel predictions: asymmetric fraying (peripheral properties lost before core ones) and sticky-lock drift (entrenched categories resist change even after their stabilizing mechanisms weaken). Explicit falsifiers. Research agenda.

\labelsc{Back matter:} Postscript, Glossary, Formalization Appendix, Predictions Appendix, Indexes -- technical resources and comprehensive indexing (subject, author, language).

\section{Market Positioning}

\subsection*{Primary Audience}

The primary readership is researchers and advanced graduate students in theoretical linguistics, grammatical theory, linguistic typology, cognitive and usage-based linguistics, and philosophy of linguistics. The book addresses foundational questions about the nature of linguistic categories -- questions that cut across subdisciplines. Syntacticians asking what a grammatical category is, semanticists asking how definiteness works, typologists asking whether cross-linguistic categories are real, and acquisitionists asking how children learn from gradient input will all find direct engagement with their empirical domains. The argument is technical but accessible; it presupposes familiarity with grammatical analysis but not with philosophy of science.

\subsection*{Secondary Audiences}

Philosophy of language and philosophy of science represent significant secondary audiences. The book offers the first systematic application of HPC theory to linguistic categories and develops original extensions (field-relative projectibility, multi-tier coupling, the projectibility-as-interpretant synthesis). Philosophers working on natural kinds, categorization, and scientific realism will find a substantive new case study with worked-out consequences.

Cognitive scientists, particularly those working on concepts, categorization, and learning, will recognize connections to prototype theory, conceptual combination, and inductive inference. The book provides a mechanistic account of how categories remain stable despite probabilistic variation and gradient judgments -- a question central to cognitive science but rarely addressed with this level of linguistic detail.

Sociolinguists and sign language researchers form a further audience, given the book's treatment of social stabilization (Ch 16) and modality independence (Ch 13).

\subsection*{Comparable and Adjacent Titles}

No existing monograph combines this framework and scope, but several books are adjacent in important ways:

\begin{itemize}
  \item \labelsc{Croft 2001.} \emph{Radical Construction Grammar} -- shares skepticism about cross-linguistic categories but offers nominalism, not HPC realism.
  \item \labelsc{Bybee 2010.} \emph{Language, Usage and Cognition} -- documents usage-based mechanisms (entrenchment, frequency effects) but does not provide a metaphysical framework.
  \item \labelsc{Haspelmath 2010.} \enquote{Framework-free grammatical theory} -- argues for comparative concepts over universals; HPC theory provides the missing ontology.
  \item \labelsc{Khalidi 2013.} \emph{Natural Categories and Human Kinds} -- philosophy of science foundation but no linguistic application.
  \item \labelsc{Goldberg 2019.} \emph{Explain Me This} -- acquisition and generalization mechanisms but not a metaphysical account of categories as kinds.
\end{itemize}

The manuscript synthesizes these lines of work into a unified framework with systematic linguistic application.

\subsection*{Course Adoption Potential}

The book is not a textbook, but it will be assigned in graduate seminars on grammatical theory, philosophy of linguistics, usage-based approaches, and linguistic typology. It addresses questions that emerge naturally in advanced syntax and semantics courses (What is a grammatical category? Why do diagnostics conflict? Are categories discrete or gradient?) and provides a principled answer with worked examples. The final chapter's research agenda makes it especially valuable for seminars preparing students for dissertation research.

Upper-level courses in linguistics and cognitive science may assign selected chapters -- especially Chs 1--3, 9--10, and 14--15 -- as accessible introductions to foundational debates.

\subsection*{Why Now?}

Four converging developments make this book timely:

\begin{enumerate}
  \item \labelsc{Empirical pressure.} Evidence from gradient acceptability, corpus frequency effects, cross-linguistic variation, and acquisitional findings has accumulated to the point where essentialist foundations are untenable. The field has the data but not the framework.
  \item \labelsc{Theoretical readiness.} Usage-based, constructionist, and sociolinguistic traditions have identified the relevant mechanisms (frequency, entrenchment, alignment, transmission) but lack a unifying metaphysical framework. HPC theory provides it.
  \item \labelsc{Interdisciplinary momentum.} Philosophy of science has developed HPC theory substantially since Boyd's original work (Khalidi, Slater, Magnus, Craver); cognitive science has moved toward mechanistic explanation; linguistics is ready for the synthesis.
  \item \labelsc{External urgency.} The rapid expansion of large language models has foregrounded questions about what grammatical categories are and whether they transfer across languages -- questions this book addresses from first principles.
\end{enumerate}

The manuscript arrives at the moment when the field has the empirical base and theoretical tools to take this step but hasn't yet done so.

\section{Suggested Reviewers}

\noindent\labelsc{William Croft} (Professor Emeritus, University of New Mexico)\\
Expert in syntactic categories, typology, construction grammar, and language change. His own work on evolutionary models of language and population thinking makes him exceptionally well placed to assess the manuscript's critique of essentialism and its mechanistic alternative.

\medskip

\noindent\labelsc{Martin Haspelmath} (Max Planck Institute for Evolutionary Anthropology / University of Leipzig)\\
Leading typologist and critic of cross-linguistic categories; developed the comparative concepts framework. Would provide rigorous assessment of whether HPC theory answers typological objections to universals.

\medskip

\noindent\labelsc{Muhammad Ali Khalidi} (Presidential Professor, CUNY Graduate Center)\\
Philosopher of science specializing in natural kinds, cognitive categories, and HPC theory. The ideal reviewer to evaluate the philosophical rigour of the framework's translation into linguistics and the original extensions (field-relative projectibility, multi-tier coupling).

\medskip

\noindent\labelsc{Adele Goldberg} (Princeton University)\\
Leading construction grammarian; work on argument structure and category learning. Would assess whether the framework successfully integrates constructionist insights and whether it has consequences for acquisition theory.

\medskip

\noindent\labelsc{Stefan Th.\ Gries} (University of California, Santa Barbara)\\
Expert in corpus linguistics, quantitative methods, and cognitive linguistics. Would provide rigorous evaluation of the book's treatment of gradience, clustering, and empirical diagnostics.

\medskip

\noindent\labelsc{Peter Harder} (University of Copenhagen)\\
Research on functional linguistics, semantics, philosophy of language, and the social-cognitive character of linguistic norms. Especially appropriate for the book's arguments about projectibility and inter-speaker stabilization.

\end{document}
